In this section, allow us to define a few lemmas, definitions, and conventions
that serve as foundation for main results of this manuscript.
For multifold sums of powers, we use the notation provided by Donald Knuth in~\cite{knuth1993johann},
because it is simple and beautiful
\begin{definition} [Multifold sums of powers]
  \label{def:multifold-sums-of-powers-recurrence}
  For non-negative integers $r,n,m$
  \begin{align*}
    \KnuthRFoldSum{0}{n}{m}   &= n^m, \\
    \KnuthRFoldSum{1}{n}{m}   &= \KnuthRFoldSum{0}{1}{m}
    + \KnuthRFoldSum{0}{2}{m} + \cdots + \KnuthRFoldSum{0}{n}{m}, \\
    \KnuthRFoldSum{r+1}{n}{m} &= \KnuthRFoldSum{r}{1}{m}
    + \KnuthRFoldSum{r}{2}{m} + \cdots + \KnuthRFoldSum{r}{n}{m}.
  \end{align*}
\end{definition}

We utilize the following notation for falling factorials
\begin{definition} [Falling factorials] For integers $x,n$
  \begin{align*}
    \fallingFactorial{n}{k} =
    \begin{cases}
      1, & \mathrm{if \; } k = 0, \\
      n(n-1)(n-2)\cdots(n-k+1) = \prod_{j=0}^{k-1} (n-j), & \mathrm{if \; } k > 0, \\
      0, & \mathrm{else}.
    \end{cases}
  \end{align*}
\end{definition}
We utilize the following notation for central factorials
\begin{definition}[Central factorials]
  \label{def:central-factorials}
  For integers $n,k$
  \begin{align*}
    \centralFactorial{n}{k} =
    \begin{cases}
      0, \quad & \mathrm{if \; } k < 0, \\
      1, \quad & \mathrm{if \; } k = 0, \\
      n \left( n + \frac{k}{2} -1 \right)\left( n + \frac{k}{2} -2 \right)
      \cdots \left( n - \frac{k}{2} +1 \right)
      = n \fallingFactorial{n+\frac{k}{2}-1}{k-1}, \quad & \mathrm{if \; } k>0,
    \end{cases}
  \end{align*}
  where
  $\fallingFactorial{n+\frac{k}{2}-1}{k-1}
  = \prod_{j=0}^{k-2} \left( n+\frac{k}{2}-1 -j \right)$
  are falling factorials.
\end{definition}
We encourage the audience to read
J. F. Steffensen's work on the definition of
generalized factorials~\cite{steffensen1933definition}.

\begin{lemma}[Central Newton's formula]
  For a function $f$ defined on integers and arbitrary integers $x,t$
  \label{lem:central-newtons-formula}
  \begin{align*}
    f(x)
    = \sum_{k=0}^{\infty} \frac{\centralFactorial{(x-t)}{k}}{k!} \delta^{k} f(t),
  \end{align*}
  where $\delta^{k} f(t)$ denotes the $k$-th central finite difference
  \begin{align*}
    \delta^{k} f(t) = \tsum_{j=0}^{k} (-1)^j \tbinom{k}{j} f\!\left(t+\tfrac{k}{2} - j\right),
  \end{align*}
  and $\centralFactorial{(x-t)}{k}$ are central factorials
  defined by~\eqref{def:central-factorials}.
  \begin{proof}
    See~\cite[p.~462]{butzer1989central}, and~\cite[section 2, eq. (19)]{steffensen1927interpolation},
    and~\cite[section 6, eq. (21)]{riordan1968combinatorial}.
  \end{proof}
\end{lemma}
We have
\begin{lemma}[Newton's formula for powers]
  \label{lem:central-newtons-formula-for-powers}
  For non-negative integers $n,m$, and an arbitrary integer $t$,
  \begin{align*}
    n^m = \sum_{k=0}^{m} \frac{\centralFactorial{(n-t)}{k}}{k!} \delta^{k} t^m.
  \end{align*}
  \begin{proof}
    Special case of Lemma~\eqref{lem:central-newtons-formula} for $f(n) = n^m$.
  \end{proof}
\end{lemma}
We have
\begin{lemma}[Central factorials binomial form]
  \label{lem:central-factorials-binomial-form}
  For integers $n$ and $k\neq0$,
  \begin{align*}
    \frac{\centralFactorial{n}{k}}{k!}
    = \frac{n}{k} \binom{n+\frac{k}{2}-1}{k-1}.
  \end{align*}
  \begin{proof}
    We have
    \begin{align*}
      \frac{\centralFactorial{n}{k}}{k!}
      &= \frac{n}{k!}\fallingFactorial{n+\frac{k}{2}-1}{k-1}
      = \frac{n}{k (k-1)!} \fallingFactorial{n+\frac{k}{2}-1}{k-1}
      = \frac{n}{k} \binom{n+\frac{k}{2}-1}{k-1},
    \end{align*}
    because of the identity in falling factorials
    $\frac{\fallingFactorial{x}{n}}{n!} = \binom{x}{n}$.
  \end{proof}
\end{lemma}
% We have
% \begin{lemma}[Central factorials binomial form decomposition]
%   \label{lem:central-factorials-binomial-form-decomposition}
%   For  non-negative integers $n,k$,
%   \begin{align*}
%     \frac{\centralFactorial{n}{k}}{k!}
%     = \frac{1}{2}
%     \left[ \binom{n+\frac{k}{2}}{k} + \binom{n+\frac{k}{2}-1}{k} \right].
%   \end{align*}
%   \begin{proof}
%     We have $\binom{a}{k} = \frac{a}{k} \binom{a-1}{k-1}$.
%     Let $a=n + \frac{k}{2}$, then
%     \begin{align*}
%       \tbinom{n + \frac{k}{2}}{k} = \tfrac{n + \tfrac{k}{2}}{k} \tbinom{n+\frac{k}{2}-1}{k-1}
%       = \tfrac{n}{k} \tbinom{n+\frac{k}{2}-1}{k-1} + \tfrac{1}{2} \tbinom{n+\frac{k}{2}-1}{k-1}.
%     \end{align*}
%     Thus,
%     \begin{align*}
%       \tfrac{n}{k} \tbinom{n+\frac{k}{2}-1}{k-1} = \tbinom{n + \frac{k}{2}}{k} - \tfrac{1}{2} \tbinom{n+\frac{k}{2}-1}{k-1}.
%     \end{align*}
%     By moving under common denominator, and applying binomial recurrence yields
%     \begin{align*}
%       \tfrac{n}{k} \tbinom{n+\frac{k}{2}-1}{k-1}
%       = \tfrac{1}{2}
%       \left[
%         2 \tbinom{n + \frac{k}{2}}{k} - \tbinom{n+\frac{k}{2}-1}{k-1}
%       \right]
%       = \tfrac{1}{2}
%       \left[
%         \tbinom{n + \frac{k}{2}}{k} + \tbinom{n + \frac{k}{2}-1}{k} + \tbinom{n + \frac{k}{2}-1}{k-1} - \tbinom{n+\frac{k}{2}-1}{k-1}
%       \right].
%     \end{align*}
%     Thus, the statement follows.
%   \end{proof}
% \end{lemma}
We have
\begin{lemma}[Generalized hockey-stick identity]
  \label{lem:generalized-hockey-stick-identity}
  For arbitrary integers $a,b$ and $j$,
  \begin{align*}
    \sum_{k=a}^{b} \binom{k}{j} = \binom{b+1}{j+1} - \binom{a}{j+1}.
  \end{align*}
  \begin{proof}
    We have,
    $\sum_{k=a}^{b} \binom{k}{j}
    = \binom{a}{j} + \binom{a+1}{j} + \cdots + \binom{b}{j}
    = \left( \sum_{k=0}^{b} \binom{k}{j} \right) - \left( \sum_{k=0}^{a-1} \binom{k}{j} \right)$.
    By hockey-stick identity $\sum_{k=0}^{n} \binom{k}{j} = \binom{n+1}{j+1}$ yields,
    \begin{align*}
      \tsum_{k=a}^{b} \tbinom{k}{j}
      = \left( \tsum_{k=0}^{b} \tbinom{k}{j} \right) - \left( \tsum_{k=0}^{a-1} \tbinom{k}{j} \right)
      = \tbinom{b+1}{j+1} - \tbinom{a}{j+1}.
    \end{align*}
    This completes the proof.
  \end{proof}
\end{lemma}
We have
\begin{lemma}[Centered hockey-stick identity]
  \label{lem:centered-hockey-stick-identity}
  For integer $n\geq1$, and arbitrary integers $t,k,r$
  \begin{align*}
    \sum_{j=1}^{n} \binom{j -t + \frac{k}{2} +r}{k}
    = \binom{n -t + \frac{k}{2} + r +1}{k+1} - \binom{1 -t + \tfrac{k}{2} +r}{k+1}.
  \end{align*}
  \begin{proof}
    By setting $a=1 -t + \tfrac{k}{2} +r$, and $b=n-t+ \tfrac{k}{2} +r$
    into Lemma~\eqref{lem:generalized-hockey-stick-identity} yields
    \begin{align*}
      \tsum_{j=1}^{n} \tbinom{j -t + \tfrac{k}{2} +r}{k} =
      \tsum_{j=1 -t + \tfrac{k}{2} +r}^{n -t + \tfrac{k}{2} +r} \tbinom{j}{k}
      = \tbinom{n -t + \tfrac{k}{2} + r +1}{k+1} - \tbinom{1 -t + \tfrac{k}{2} +r}{k+1}.
    \end{align*}
    This completes the proof.
  \end{proof}
\end{lemma}
We have
\begin{lemma}[Multifold sums of zero powers]
  \label{lem:multifold-sum-of-zero-powers}
  For non-negative integers $r,n$
  \begin{align*}
    \KnuthRFoldSum{r}{n}{0} = \binom{r+n-1}{r}.
  \end{align*}
  \begin{proof}
    \leavevmode\par
    \begin{enumerate}
      \item Let $r=0$, then we have $\KnuthRFoldSum{0}{n}{0} = 1$,
        by definition~\eqref{def:multifold-sums-of-powers-recurrence}.
      \item Let $r=1$, then we have
        $\KnuthRFoldSum{1}{n}{0}
        = \sum_{k=1}^{n} 1 = \binom{n}{1}$.
      \item Let $r=2$, then we have
        $\KnuthRFoldSum{2}{n}{0}
        = \sum_{k=1}^{n} \binom{k}{1}
        = \binom{n+1}{2}$.
      \item Let $r=3$, then we have
        $\KnuthRFoldSum{3}{n}{0}
        = \sum_{k=1}^{n} \binom{k+1}{2}
        = \binom{n+2}{3}$.
      \item Similarly, for arbitrary integer $r$, by hockey-stick identity
        $\sum_{k=1}^{n} \binom{k}{r} = \binom{n+1}{r+1}$,
        yields
        $\KnuthRFoldSum{r}{n}{0}
        =  \sum_{k=1}^{n} \KnuthRFoldSum{r-1}{k}{0}
        =  \sum_{k=1}^{n} \binom{k+r-2}{r-1}
        = \binom{r+n-1}{r}$.
    \end{enumerate}
    This completes the proof.
  \end{proof}
\end{lemma}
We follow the convention
\begin{convention} For all $x$
  \begin{align*}
    x^0 = 1.
  \end{align*}
\end{convention}
Donald Knuth give extensive discussion on the convention $x^0=1$ for
all $x$, including zero, in Concrete Mathematics~\cite[~p. 162]{graham1994concrete},
and in~\cite{knuth1992notesnotation}.
